\section{Experimental Setup}
\label{sec:experiments}

We design experiments to test each step of the hypothesized causal chain, from counterfactual information extraction through targeted editing to fix success.

\subsection{Research Questions}

Our experiments address four questions, each mapping to a specific claim:

\textbf{RQ1 (H-M1):} Do execution traces contain extractable counterfactual information?\\
\emph{Tests:} Whether error traces provide the localization information our theory predicts.

\textbf{RQ2 (H-M2):} Does detailed feedback enable targeted edits compared to binary feedback?\\
\emph{Tests:} Whether information granularity affects model editing behavior.

\textbf{RQ3 (H-M3):} Do targeted edits achieve higher fix rates than global rewrites?\\
\emph{Tests:} Whether the predicted mechanism (targeting $\rightarrow$ success) holds empirically.

\textbf{RQ4 (H-C1):} Does execution advantage increase with task complexity?\\
\emph{Tests:} Whether complex tasks benefit more from precise localization.

\subsection{Datasets}

We evaluate on two standard code generation benchmarks with complementary characteristics:

\begin{table}[h]
\centering
\begin{tabular}{lllll}
\toprule
\textbf{Dataset} & \textbf{Problems} & \textbf{Language} & \textbf{Complexity} & \textbf{Why Chosen} \\
\midrule
HumanEval & 164 & Python & Lower & Standard benchmark; comparison with prior work \\
MBPP & 500 & Python & Higher & Tests complexity interaction hypothesis \\
\bottomrule
\end{tabular}
\end{table}

\textbf{HumanEval}~\cite{chen2021humaneval} provides function-level code completion tasks with doctests. We use the standard 164-problem split. Test cases are deterministic, enabling clean execution feedback collection.

\textbf{MBPP}~\cite{austin2021mbpp} contains more complex programming problems with multi-step reasoning. We use this to test whether execution advantage increases with task complexity (RQ4).

\subsection{Feedback Conditions}

We compare four feedback conditions:

\begin{table}[h]
\centering
\begin{tabular}{lll}
\toprule
\textbf{Condition} & \textbf{Information Content} & \textbf{Purpose} \\
\midrule
Execution-Detailed & Line number, error type, expected/actual & Full localization \\
Execution-Binary & Pass/fail only & Ablates localization \\
AI-Critic & LLM-generated critique & Tests AI approximation \\
Random & Shuffled feedback & Controls for feedback presence \\
\bottomrule
\end{tabular}
\end{table}

The random condition provides a critical control: if AI-critic performs at random level, the critic provides no useful signal. If AI-critic exceeds random but underperforms execution, we can quantify the localization gap.

\subsection{Model and Training}

\textbf{Base Model:} CodeLlama-7B-Instruct (Meta, 2023)

We select this model for three reasons: (1) open-source, enabling reproducibility; (2) instruction-tuned for code editing; (3) representative of widely-deployed 7B code LLMs.

\textbf{Refinement Protocol:}
\begin{itemize}
\item Iterations: $k=3$
\item Temperature: 0.2 (low variance for reproducibility)
\item Max tokens: 512 per generation
\item Timeout: 10s per execution
\end{itemize}

\textbf{AI-Critic:} CodeLlama-7B-Instruct as critic (same model, different prompt). This controls for model capability---any execution advantage cannot be attributed to using a stronger critic.

\subsection{Evaluation Metrics}

\textbf{Primary:}
\begin{itemize}
\item \textbf{pass@1:} Proportion of problems solved after $k=3$ refinement iterations. Computed per condition.
\end{itemize}

\textbf{Mechanism Metrics:}
\begin{itemize}
\item \textbf{CF-score:} Counterfactual information score (0--1) per trace. Measures line number, error type, expected/actual value presence.
\item \textbf{Edit scope:} Classified as Targeted ($\leq$5 lines), Local (6--15 lines), or Global ($>$15 lines).
\item \textbf{Fix rate by scope:} Success rate conditioned on edit scope.
\end{itemize}

\textbf{Statistical Testing:}
\begin{itemize}
\item Paired t-test for within-condition comparisons
\item Bonferroni correction for multiple comparisons
\item Effect size reported as Cohen's $d$
\end{itemize}

\subsection{Implementation Details}

\textbf{Execution Sandbox:}
\begin{itemize}
\item Docker containers with Python 3.10
\item pytest with captured stdout/stderr
\item 10-second timeout per execution
\item 512MB memory limit
\end{itemize}

\textbf{Trace Parsing:}
\begin{itemize}
\item Regex extraction of line numbers, error types
\item Template-based conversion to natural language
\item CF-score computed as weighted sum of present elements
\end{itemize}

\textbf{Edit Analysis:}
\begin{itemize}
\item Line-level diff (difflib)
\item AST-based change localization
\item Automatic scope classification
\end{itemize}

\textbf{Compute:}
\begin{itemize}
\item Hardware: 1$\times$ NVIDIA A100 (40GB)
\item Inference time: $\sim$0.5s per generation
\item Total experiment runtime: $\sim$48 hours
\end{itemize}

All experiments use fixed random seeds for reproducibility. Code and data will be released upon publication.
